
\documentclass{article}
\usepackage{colortbl}
\usepackage{makecell}
\usepackage{multirow}
\usepackage{supertabular}

\begin{document}

\newcounter{utterance}

\centering \large Interaction Transcript for game `referencegame', experiment `letter\_grids', episode 11 with reflexion{-}Qwen3.5{-}4B.
\vspace{24pt}

{ \footnotesize  \setcounter{utterance}{1}
\setlength{\tabcolsep}{0pt}
\begin{supertabular}{c@{$\;$}|p{.15\linewidth}@{}p{.15\linewidth}p{.15\linewidth}p{.15\linewidth}p{.15\linewidth}p{.15\linewidth}}
    \# & $\;$A & \multicolumn{4}{c}{Game Master} & $\;\:$B\\
    \hline

    \theutterance \stepcounter{utterance}  
    & & \multicolumn{4}{p{0.6\linewidth}}{
        \cellcolor[rgb]{0.9,0.9,0.9}{
            \makecell[{{p{\linewidth}}}]{
                \texttt{\tiny{[P1$\langle$GM]}}
                \texttt{You have attempted this same task before and did not succeed. These are the plans you wrote after those attempts.} \\
\texttt{Trial \#1: New plan: My previous description was too vague because it failed to uniquely distinguish the target grid from the distractors. Specifically, I missed the critical difference in the third row (target has one 'X' at the end, distractors have none) and the fourth/fifth rows (target has 'X's in the first column, while distractors have 'X's in the first column but different patterns in the bottom rows). Next time, I will construct the expression by explicitly listing the exact coordinates or patterns of the 'X's that differ from the distractors, ensuring every unique feature of the target grid is mentioned to guarantee uniqueness. For example, I will specify: "The grid with a full top row, a full first column, a third row with only the last cell filled, and a fourth row with only the first cell filled."} \\
\\ 
\\ 
\texttt{You are given three grids, where each of them is 5 by 5 in size.} \\
\texttt{Grids have empty cells marked with "▢" and filled cells marked with "X".} \\
\texttt{Your task is to generate a referring expression that best describes the target grid while distinguishing it from the two other distractor grids.} \\
\texttt{The first grid is the target grid, and the following two grids are the distractors.} \\
\\ 
\texttt{Target grid:} \\
\\ 
\texttt{X X X X X} \\
\texttt{□ □ □ □ □} \\
\texttt{X X X X □} \\
\texttt{X □ □ □ □} \\
\texttt{X □ □ □ □} \\
\\ 
\texttt{Distractor grid 1:} \\
\\ 
\texttt{X X X X X} \\
\texttt{X □ □ □ □} \\
\texttt{X X X X X} \\
\texttt{X □ □ □ □} \\
\texttt{X X X X X} \\
\\ 
\texttt{Distractor grid 2:} \\
\\ 
\texttt{X X X X X} \\
\texttt{X □ □ □ □} \\
\texttt{X □ □ □ □} \\
\texttt{X □ □ □ □} \\
\texttt{X X X X X} \\
\\ 
\texttt{Instruction: Describe the target grid.} \\
\texttt{Generate the referring expression starting with the tag "Expression: " for the given target grid. Omit any other text.} \\
            }
        }
    }
    & \\ \\

    \theutterance \stepcounter{utterance}  
    & \multicolumn{4}{p{0.6\linewidth}}{
        \cellcolor[rgb]{0.9,0.9,0.9}{
            \makecell[{{p{\linewidth}}}]{
                \texttt{\tiny{[P1$\rangle$GM]}}
                \texttt{Expression: The grid with a full top row, a full first column, a third row with only the last cell filled, and a fourth row with only the first cell filled.} \\
            }
        }
    }
    & & \\ \\

    \theutterance \stepcounter{utterance}  
    & & & \multicolumn{2}{p{0.3\linewidth}}{
        \cellcolor[rgb]{0.9,0.9,0.9}{
            \makecell[{{p{\linewidth}}}]{
                \texttt{\tiny{[GM$|$GM]}}
                \texttt{Expression: The grid with a full top row, a full first column, a third row with only the last cell filled, and a fourth row with only the first cell filled.} \\
            }
        }
    }
    & & \\ \\

    \theutterance \stepcounter{utterance}  
    & & \multicolumn{4}{p{0.6\linewidth}}{
        \cellcolor[rgb]{0.9,0.9,0.9}{
            \makecell[{{p{\linewidth}}}]{
                \texttt{\tiny{[GM$\rangle$P2]}}
                \texttt{You have attempted this same task before and did not succeed. These are the plans you wrote after those attempts.} \\
\texttt{Trial \#1: New plan: I incorrectly identified the third row of the "Second" grid as having a full row except for the last cell, when in fact it was completely empty. I failed to carefully verify the specific pattern of filled ("X") and empty ("▢") cells in each row against the description. Next time, I will systematically check each row of every grid individually: first confirming the top row is full, then verifying the third row has exactly one empty cell at the end, and finally ensuring the first column has only one empty cell at the second position. I will only select a grid if all three conditions match perfectly before answering.} \\
\\ 
\\ 
\texttt{You are given three grids, where each of them is 5 by 5 in size.} \\
\texttt{Grids have empty cells marked with "▢" and filled cells marked with "X".} \\
\texttt{You are also given a referring expression that describes one of the given grids.} \\
\texttt{Your task is to select the grid that matches the given referring expression.} \\
\texttt{Generate only the number (in text) of the grid that matches the given expression by selecting first, second, or third.} \\
\\ 
\texttt{First:} \\
\\ 
\texttt{X X X X X} \\
\texttt{X □ □ □ □} \\
\texttt{X □ □ □ □} \\
\texttt{X □ □ □ □} \\
\texttt{X X X X X} \\
\\ 
\texttt{Second:} \\
\\ 
\texttt{X X X X X} \\
\texttt{X □ □ □ □} \\
\texttt{X X X X X} \\
\texttt{X □ □ □ □} \\
\texttt{X X X X X} \\
\\ 
\texttt{Third:} \\
\\ 
\texttt{X X X X X} \\
\texttt{□ □ □ □ □} \\
\texttt{X X X X □} \\
\texttt{X □ □ □ □} \\
\texttt{X □ □ □ □} \\
\\ 
\texttt{Expression: The grid with a full top row, a full first column, a third row with only the last cell filled, and a fourth row with only the first cell filled.} \\
\texttt{Question: Which grid does the expression refer to?} \\
\texttt{Start with the tag "Answer: ", followed by your selection. Omit any other text.} \\
            }
        }
    }
    & \\ \\

    \theutterance \stepcounter{utterance}  
    & & & \multicolumn{4}{p{0.6\linewidth}}{
        \cellcolor[rgb]{0.9,0.9,0.9}{
            \makecell[{{p{\linewidth}}}]{
                \texttt{\tiny{[GM$\langle$P2]}}
                \texttt{Answer: First} \\
            }
        }
    }
     \\ \\

    \theutterance \stepcounter{utterance}  
    & & & \multicolumn{2}{p{0.3\linewidth}}{
        \cellcolor[rgb]{0.9,0.9,0.9}{
            \makecell[{{p{\linewidth}}}]{
                \texttt{\tiny{[GM$|$GM]}}
                \texttt{first} \\
            }
        }
    }
    & & \\ \\

    \theutterance \stepcounter{utterance}  
    & \multicolumn{4}{p{0.6\linewidth}}{
        \cellcolor[rgb]{0.95,0.95,0.95}{
            \makecell[{{p{\linewidth}}}]{
                \texttt{\tiny{[P2$\rightarrow$P1]}}
                \texttt{Post{-}episode Reflexion (not sent to GM):} \\
\texttt{New plan: My previous description failed because it was factually incorrect regarding the target grid's structure; I claimed the fourth row had only the first cell filled, but the target grid actually has the first cell filled in the fourth and fifth rows, while the third row has the last cell filled. To succeed, I must accurately describe the specific pattern of 'X's in the target grid relative to the distractors, ensuring I correctly identify that the target has a full top row, a full first column, a third row with only the last cell filled, and a fourth row with only the first cell filled (and the fifth row also with only the first cell filled, which distinguishes it from distractor 2). I will generate the expression by explicitly stating these precise row patterns to ensure uniqueness.} \\
            }
        }
    }
    & & \\ \\

    \theutterance \stepcounter{utterance}  
    & & & \multicolumn{4}{p{0.6\linewidth}}{
        \cellcolor[rgb]{0.95,0.95,0.95}{
            \makecell[{{p{\linewidth}}}]{
                \texttt{\tiny{[P1$\rightarrow$P2]}}
                \texttt{Post{-}episode Reflexion (not sent to GM):} \\
\texttt{New plan: I misidentified the pattern of the third row in the "Second" grid, assuming it had a filled cell at the end when it was actually empty. I will now strictly verify each grid against the expression's four specific conditions: 1) Top row is all "X", 2) First column is all "X", 3) Third row has only the last cell as "X" (others "▢"), and 4) Fourth row has only the first cell as "X" (others "▢"). I will compare these exact patterns against the provided grids before selecting the correct number.} \\
            }
        }
    }
     \\ \\

\end{supertabular}
}

\end{document}
